\section{Line Integrals}

\subsection{Mistakes and Tips}

\begin{mistake}[What actualy is \(f(\mathbf{x}(t))\)?]

\(f(\mathbf{x}(t))\) is a function that acts on the entries for \(\mathbf{x}(t) = (x_1(t),x_2(t),x_3(t))\) of a curve. 

\end{mistake}

\begin{mistake}[\(\mid \frac{d \mathbf{x}}{dt} \mid \) differs]
  \(\mid \frac{d \mathbf{x}}{dt} \mid \) is not always a fixed number and may be a function of \(t\) that differs.
\end{mistake}

\begin{mistake}[Parametrising implicit curves]
  Sometimes you need to parametrise the function before you can integrate under the scalar field. Make sure to remember your polar coordinates and \textb{draw a diagran} if possible
\end{mistake}

\begin{mistake}[Splitting piecewise smooth curves]
  When a curve is piecwise smooth, make sure to split it when integrating.
\end{mistake}

\begin{mistake}[Unit tangents and orthonormal]
  You can always find the unit tangent vectors to a curve but they are not always orthonormal
\end{mistake}

\subsection{Methods}

\begin{nmeth}[Evaluating a Scalar Line Integral]\emph{}
  \begin{enumerate}
    \item Write $f(\mathbf{x}), \mathbf{x}(t) = (x_1(t),x_2(t),x_3(t))$ and $t \in [t_0,t_1]$
    \item Find $\frac{d \mathbf{x}}{dt}$ and then take the magnitude (this may not be an integer but a function of $t$)
    \item Substitute in and integrate
  \end{enumerate}
\end{nmeth}


\begin{nmeth}[Evaluating arc length]\emph{}
  \begin{enumerate}
    \item Write $\mathbf{x}(t) = (x_1(t),x_2(t),x_3(t))$
    \item Find $\frac{d \mathbf{x}}{dt}$
    \item Perform $\sqrt{\frac{d \mathbf{x}}{dt} \cdot \frac{d \mathbf{x}}{dt}}$
    \item Use the formula for arclength
  \end{enumerate}
\end{nmeth}

\begin{nmeth}[Evaluating a Vector Field Line Integral]\emph{}
  
\begin{enumerate}
  \item Sketch $\mathbf{x}(t)$ if possible
  \item Substitute $\mathbf{x}(t)$ into $\mathbf{F(x)}$ to form $\mathbf{F(x(}t))$
  \item Find $\frac{d\mathbf{x}}{dt}$
  \item Perform $\mathbf{F(x(}t)) \cdot \frac{d \mathbf{x}}{dt}$
  \item Substitute into the formula 
\end{enumerate}

\end{nmeth}

\begin{nmeth}[Finding unit tangent vectors]
  To find the unit tangent vectors for the coordinate system $\mathbf{x}(u,v,w)$ use the formulas:
  $$
\boldsymbol{e}_u=\frac{1}{h_u} \frac{\partial \boldsymbol{x}}{\partial u}, \quad \boldsymbol{e}_v=\frac{1}{h_v} \frac{\partial \boldsymbol{x}}{\partial v}, \quad \boldsymbol{e}_w=\frac{1}{h_w} \frac{\partial \boldsymbol{x}}{\partial w},
$$

where

$$
h_u=\left|\frac{\partial \boldsymbol{x}}{\partial u}\right|, \quad h_v=\left|\frac{\partial \boldsymbol{x}}{\partial v}\right|, \quad h_w=\left|\frac{\partial \boldsymbol{x}}{\partial w}\right|
$$

are called the scale factors.
\end{nmeth}

\begin{nmeth}[Ensuring an Orthonormal Basis]
  To check if these unit tangent vectors form an orthonormal basis we will have to ensure that they are all mutually orthogonal.
\end{nmeth}
